% 2-5-hybrid.tex  | Budget 2pp
% SOURCE MAP: hybrid DNN ->kuo ; 4-module hybrid (isolation forest+error comp+prob)
%  ->zhang_deep_2020 ; ARD+bagging ExtraTrees ->alkawaz ; ML+bionic optim ->sun ;
%  MIMO point+interval ->jiang_multivariable_2023 ; multi-head attention+CNN
%  ->pourdaryaei ; SDAE+BiLSTM+GRU ->moradzadeh ; CNN-BiLSTM+AR ->mubarak ;
%  ensemble ML ->shah ; adaptive VMD+kernel ->yang ; 2-step DL+err-comp AU ->ghimire.
% Eq = weighted ensemble combination. ghimire abstract-less: described from title only.
\section{مدل‌های ترکیبی و مکانیزم توجه}

محدودیت‌های مدل‌های منفرد \lr{—} چه آماری و چه یادگیری عمیق \lr{—} پژوهشگران را
به سوی مدل‌های ترکیبی سوق داده است؛ رویکردی که می‌کوشد نقاطِ قوتِ چند روش را
کنارِ هم بگذارد. منطقِ بنیادیِ بسیاری از این روش‌ها، ترکیبِ وزن‌دارِ خروجیِ چند
مدلِ پایه است:

\begin{equation}
\hat{y} = \sum_{k=1}^{K} w_k\,\hat{y}_k
\end{equation}

که در آن $\hat{y}_k$ پیش‌بینیِ مدلِ $k$‌اُم، $w_k$ وزنِ آن، و $\hat{y}$ پیش‌بینیِ
نهاییِ ترکیب است. این ایده در ادبیات به شکل‌های گوناگون تحقق یافته است؛ از ترکیبِ
رگرسیونِ خطی و درخت‌های تصادفیِ کیسه‌ای \cite{alkawaz_day-ahead_2022} تا مدل‌های
گروهیِ متشکل از روش‌های آماری و شبکه‌های عصبی \cite{shah_short-term_2021}.

الگوی پرتکرارِ دیگر، ترکیبِ پیش‌پردازش و تجزیه‌ی سیگنال با یادگیری عمیق است.
رویکردهایی مبتنی بر تجزیه‌ی مُدِ متغیر و حذفِ نوفه، سریِ قیمت را پیش از مدل‌سازی
به مؤلفه‌های ساده‌تر تجزیه می‌کنند و سپس آن‌ها را به شبکه‌های بازگشتی می‌سپارند
\cite{yang_novel_2022, moradzadeh_hybrid_2025}. چارچوب‌های چندپیمانه‌ای نیز
پیشنهاد شده‌اند که در آن‌ها ماژول‌های جداگانه‌ای برای پیش‌پردازش، پیش‌بینیِ
نقطه‌ای، جبرانِ خطا و پیش‌بینیِ احتمالاتی تعبیه شده است
\cite{zhang_deep_2020, ghimire_two-step_2024}. افزودنِ مکانیزمِ توجه و شبکه‌های
کانولوشنی-بازگشتی نیز از جدیدترین گرایش‌ها در این حوزه است
\cite{pourdaryaei_new_2024, mubarak_day-ahead_2024, kuo_electricity_2018}.

سومین گرایش، بهره‌گیری از الگوریتم‌های بهینه‌سازی برای تنظیمِ ابرپارامترها و
ساختارِ مدل است \cite{sun_day-ahead_2022}، و برخی چارچوب‌ها فراتر از پیش‌بینیِ
نقطه‌ای، پیش‌بینیِ بازه‌ای و چندخروجی را نیز پوشش می‌دهند
\cite{jiang_multivariable_2023}. جمع‌بندیِ این بخش آن است که مدل‌های ترکیبی و
گروهی، به‌ویژه در بازارِ روز-پیش، عملکردی قدرتمند از خود نشان می‌دهند؛ اما این
قدرت به بهای پیچیدگیِ بیشتر، دشواریِ تفسیر و خطرِ بیش‌برازش به‌دست می‌آید. همین دو
نکته \lr{—} کارآمدیِ بالای ترکیب از یک سو و کاهشِ تفسیرپذیری از سوی دیگر \lr{—}
به‌طورِ مستقیم به دو محورِ اصلیِ این پایان‌نامه، یعنی وزن‌دهیِ گروهیِ رژیم‌محور و
تحلیلِ تفسیرپذیری، پیوند می‌خورد.
